\documentclass[11pt,a4paper]{article}
\usepackage{fontspec}
\setmainfont{DejaVu Sans}[Scale=0.90]
\usepackage{amsmath,amssymb}
\usepackage[margin=2.2cm]{geometry}
\usepackage{booktabs,array,tabularx}
\usepackage{graphicx}
\graphicspath{{../figures/}}
\usepackage{enumitem}
\usepackage{titlesec}
\usepackage{parskip}
\titlelabel{\thetitle.\quad}
\titleformat{\section}{\normalfont\large\bfseries}{\thesection.}{0.6em}{}
\titleformat{\subsection}{\normalfont\normalsize\bfseries}{\thesubsection}{0.6em}{}
\titlespacing*{\section}{0pt}{15pt}{5pt}
\titlespacing*{\subsection}{0pt}{11pt}{4pt}
\setlist[itemize]{leftmargin=1.4em,itemsep=1pt,topsep=3pt}
% Preserve fonts/margins; permit a final line-breaking pass for long prose.
\setlength{\emergencystretch}{2em}
\begin{document}
\begin{center}
{\Large\bfseries Orientation et comptage des canaux cubiques}\\[4pt]
{\normalsize Du tenseur complexe à un événement de population conditionnel}\\[2pt]
{\small Partie I --- 27 septembre 2026}
\end{center}

\section{Question, convention et portée}
Les notes 20 et 21 reconstruisent des interactions et des amplitudes complexes
pour un triplet de l'AlN. Cette note précise leur traduction en canaux.
Elle ne construit ni l'opérateur de collision complet du matériau ni une
trajectoire hydrodynamique en température. Les identités de comptage sont
connues ; aucune nouveauté théorique n'est revendiquée.

Pour des fréquences strictement positives, adoptons une base réciproque
appariée, $X_\lambda=a_\lambda+a^\dagger_{\bar\lambda}$, et
\[
 H_3=\sum_{\lambda\mu\nu\ {\rm ordonnés}}
 V_{\lambda\mu\nu}X_\lambda X_\mu X_\nu,\qquad
 V_{\bar\lambda\bar\mu\bar\nu}=V_{\lambda\mu\nu}^{*}.
\]
Le tenseur est supposé symétrique sous permutation et contient déjà
$1/3!$. La convention alternative $H_3=(1/6)\sum TXXX$ donne $T=6V$.
Une base numérique arbitraire exige un raccordement des vecteurs propres
entre $q$ et $-q$ ; le simple changement d'un indice ne suffit pas.

\section{Coefficient hamiltonien et éléments de Fock}
La désintégration $p\to a+b$ utilise le triplet entrant
$(p,\bar a,\bar b)$ avec $q_p-q_a-q_b=G$.
La somme des tuples ordonnés distincts donne
\[
 H_e=g_e a_p a_a^\dagger a_b^\dagger+{\rm h.c.},
 \qquad
 g_e=\begin{cases}6V_{p\bar a\bar b},&a\ne b,\\
                   3V_{p\bar a\bar a},&a=b.\end{cases}
\]
L'égalité $a=b$ concerne le mode complet, pas seulement sa fréquence
ou son bloc dégénéré. Sans symétrie du tenseur, il faut sommer les
amplitudes complexes pertinentes avant de prendre le module au carré.

Pour les occupations entières $N$, les facteurs de création sont
\[
 |M_{pab}|^2=|g_e|^2N_p(N_a+1)(N_b+1),\qquad
 |M_{paa}|^2=|g_e|^2N_p(N_a+1)(N_a+2).
\]
Le processus inverse contient $(N_p+1)N_aN_b$, ou
$(N_p+1)N_a(N_a-1)$. Pour un parent et des filles dans le vide,
les probabilités relatives sont donc $36|V|^2$ et $18|V|^2$.
L'adjoint donne la réaction inverse sur les mêmes modes ; le canal
aux vecteurs d'onde opposés agit en général sur d'autres populations.
Il ne doit être ni supprimé ni ajouté deux fois.

\section{Hypothèse cinétique et facteur de population}
Une règle d'or nécessite un régime de couplage faible, une approximation
cinétique et une mesure spectrale justifiée. Un Hamiltonien fini résonant
ne fournit pas une valeur finie de $\delta(0)$.
La fermeture en occupations moyennes $n$ suppose en outre la
factorisation des modes. Pour une fille répétée, nous imposons les
moments d'une loi géométrique :
\[
 \langle N(N-1)\rangle=2n^2,\qquad
 \langle(N+1)(N+2)\rangle=2(1+n)^2.
\]
Cette hypothèse n'est pas une conséquence de la connaissance de $n$ :
à $n=1$, un état de Fock donne les moments $0,6$, contre $2,8$
pour la loi géométrique.

Avec $P=|V|^2$ en eV$^2$, $f$ en THz ordinaires et une densité
$D_f$ en THz$^{-1}$, l'événement réversible non ordonné s'écrit
\[
 \kappa_e=\frac{36P D_f}{\hbar_{\rm ps}^{\,2}(1+\delta_{ab})},
 \quad
 J_e=\kappa_e[n_p(1+n_a)(1+n_b)-(1+n_p)n_an_b],
 \quad \dot n=-rJ_e,
\]
où $\hbar_{\rm ps}$ est en eV\,ps, $\kappa_e$ en ps$^{-1}$ et
$r=e_p-e_a-e_b$ (donc $e_p-2e_a$ si $a=b$).
La densité doit être intégrée selon la mesure du problème ; elle n'est
pas un poids arbitraire pour un tuple désaccordé.

À l'équilibre résonant, $Q=\bar n_p(1+\bar n_a)(1+\bar n_b)$ et
$W_{ii}=\bar n_i(1+\bar n_i)$ donnent
\[
 L_e=\kappa_e Q\,rr^{\mathsf T}W^{-1},\qquad
 C_e=W^{-1/2}L_eW^{1/2}
     =\kappa_e Q\,uu^{\mathsf T},\quad u=W^{-1/2}r.
\]
Ainsi $C_e$ est positif et $C_eW^{1/2}\varepsilon=0$ si
$r^{\mathsf T}\varepsilon=0$. Ce sont des conséquences du modèle
conditionnel, sans validation de sa fermeture pour l'AlN.

Cette fermeture n'est même pas un sous-espace statistique invariant
pour le modèle de sauts de Fock déclaré. En posant $c=\kappa_e/2$,
les taux sont $cX(Y+1)(Y+2)$ et $c(X+1)Y(Y-1)$. Pour une loi
initiale géométrique produit de moyenne $(p,a)$,
\[
 \left.\frac{d}{dt}
 [\langle Y(Y-1)\rangle-2a^2]\right|_{t=0}
 =2(2a+1)J_e.
\]
Ce défaut est non nul hors équilibre et apparaît déjà au premier ordre
près de l'équilibre. La factorisation positive ci-dessus ne démontre
donc pas l'autonomie exacte des occupations moyennes.

\section{Lien précis avec la largeur de raie}
Le code phono3py 4.5.0 inclut déjà $1/36$, les trois facteurs
d'oscillateur et $1/N_q$ dans son export $P$. Il ne faut pas diviser
à nouveau par le nombre de points.
Sa demi-largeur $\gamma$ en THz donne $\tau^{-1}=4\pi\gamma$ en
ps$^{-1}$ et son facteur de conversion est
$C_\gamma=18\pi/(h\,10^{12})^2$.
Pour des filles distinctes, $\kappa_e=8\pi C_\gamma P D_f$ :
les deux ordres de filles sont présents dans la somme de raie.
Pour une fille répétée, $\kappa_e=4\pi C_\gamma P D_f$.

L'identification de chaque diagonale de population à $4\pi\gamma$
n'est cependant pas universelle. Dans la fermeture répétée ci-dessus,
la dérivée de la population fille compte ses deux occurrences :
\[
 (L_e)_{aa}=4\kappa_e(\bar n_a-\bar n_p),\qquad
 (4\pi\gamma_a)_{\rm canal}=2\kappa_e(\bar n_a-\bar n_p).
\]
Le rapport est deux pour la fille et un pour le parent. Il s'agit
d'une distinction entre les objets comparés dans ce modèle discret,
pas d'une erreur démontrée du calcul de self-énergie.

\section{Contrôles et limite matérielle}
Quatre approches indépendantes précèdent la confrontation : expansion
hamiltonienne, combinatoire des permutations, matrices de Fock, et
analyse dimensionnelle du code. Les matrices de Fock reproduisent les
coefficients avec une erreur relative inférieure à $1.5\,10^{-16}$.
Un mauvais double comptage des tuples répétés multiplie l'amplitude
par deux et la probabilité par quatre.

Pour le seul triplet AlN de lignes $(1,13,28)$ du maillage $3^3$,
les six permutations reconstruites concordent à $2.4\,10^{-15}$
en norme relative. Les lignes opposées sont $(2,26,14)$.
Le raccordement unitaire par blocs permet un contrôle de conjugaison
du canal orienté avec un résidu de $1.6\,10^{-14}$ au plus.
La liberté physique de rotation exige des fréquences égales ; les blocs
numériques et leurs petits résidus ne prouvent pas une dégénérescence exacte.

Parmi les 1472 entrées au-dessus de $10^{-12}$ du maximum de $P$,
aucune n'a $|f_0-f_1-f_2|<10^{-10}$ THz. Le désaccord minimal est
$0.0064023$ THz. Cette observation concerne cette orientation et ce seuil,
pas tous les canaux du matériau. Une largeur gaussienne non nulle
n'autorise donc pas leur import comme événements exactement résonants.

\section{Suite et sources}
La prochaine étape est une intégration d'énergie conservatrice,
avec un modèle cinétique déclaré et une énumération globale vérifiée.
Restent ensuite la convergence matérielle, la température, les frontières
du film et la confrontation expérimentale.

Sources primaires : Togo, Chaput et Tanaka, Phys.\ Rev.\ B
\textbf{91}, 094306 (2015), doi:10.1103/\allowbreak PhysRevB.91.094306 ;
Fugallo et al., Phys.\ Rev.\ B \textbf{88}, 045430 (2013),
doi:10.1103/\allowbreak PhysRevB.88.045430.
Les rapports, scripts, objections et limites sont conservés dans
\texttt{theory/aln/channel\_counting/}.
\end{document}
